% ============================================================
% CHAPITRE 1 — État des lieux du recrutement
% Structure calquée sur le rapport de référence (Ch.1 TMS)
% ============================================================
\chapter{Problèmes du Recrutement --- État des Lieux et Évolution}

\section{Introduction}

Le recrutement est devenu un pilier central de la performance organisationnelle.
Autrefois réduit à la publication d'annonces et au tri manuel de CV, ce processus
constitue désormais un levier stratégique de compétitivité, d'attractivité employeur
et de qualité des embauches. La digitalisation, la mobilité professionnelle, la
concurrence pour les talents et les attentes croissantes des candidats (réponse rapide,
transparence du parcours, expérience utilisateur fluide) renforcent cette importance.

Ainsi, le recrutement ne se limite plus à pourvoir un poste vacant~: il conditionne
la capacité d'une organisation --- en particulier d'un réseau multi-agences --- à
attirer, évaluer, sélectionner et intégrer les bons profils, tout en conservant une
traçabilité claire des décisions RH.

Ce chapitre vise à dresser un \textbf{état des lieux} du recrutement moderne, à
identifier les \textbf{limites des pratiques traditionnelles}, puis à présenter les
\textbf{tendances actuelles} des systèmes de gestion du recrutement (ATS / plateformes RH),
afin de préparer la justification de la solution développée dans le cadre de ce PFA.


\section{Situation actuelle du recrutement et des RH numériques}

Le marché du travail contemporain se caractérise par une accélération des échanges
entre employeurs et candidats. Les canaux de diffusion se sont multipliés (sites
carrière, jobboards, réseaux sociaux professionnels, cabinets, cooptation), tandis
que le volume de candidatures peut devenir difficile à absorber sans outil adapté.

Dans les organisations structurées, le processus de recrutement suit généralement
un parcours type~:
\begin{enumerate}
  \item expression du besoin et définition du poste~;
  \item publication de l'offre~;
  \item réception et tri des candidatures (CV, lettres, formulaires)~;
  \item présélection (entretiens RH, tests, QCM)~;
  \item entretiens métiers / managers~;
  \item décision (retenu, non retenu, désisté)~;
  \item confirmation d'embauche et intégration.
\end{enumerate}

Cette chaîne implique plusieurs acteurs --- RH, managers, candidats, parfois
administrateurs transverses --- et génère de nombreuses informations sensibles~:
données personnelles, documents (CV), historiques d'entretiens, scores d'évaluation,
statuts de candidature, etc.

Pour un réseau d'agences tel que \textbf{DAAM}, s'ajoute une contrainte de
\textbf{gouvernance territoriale}~: le recrutement s'organise souvent par \textbf{zones}
et \textbf{agences}, avec des responsables RH dont le périmètre doit être respecté,
tout en permettant à l'administration une \textbf{vue globale} (volume, taux de
sélection, performance par responsable, suivi des admissions).

Sur le plan technologique, les entreprises s'orientent de plus en plus vers des
solutions numériques capables de~:
\begin{itemize}
  \item centraliser les offres et les candidatures~;
  \item standardiser le workflow de sélection~;
  \item automatiser les notifications (convocation à entretien, confirmation d'embauche)~;
  \item produire des indicateurs de pilotage (KPI) utiles à la décision.
\end{itemize}

Ces besoins ouvrent la voie aux plateformes ATS (\textit{Applicant Tracking System})
et aux modules SI RH spécialisés dans le recrutement.


\section{Limites de la gestion traditionnelle du recrutement}

Malgré la digitalisation croissante, de nombreuses structures continuent de s'appuyer
sur des pratiques fragmentées~: tableurs Excel, messagerie électronique, dossiers
partagés, outils disparates (jobboard, calendrier, formulaires). Cette gestion
\og{}traditionnelle\fg{} présente plusieurs limites majeures.

\paragraph{Dispersion de l'information.}
Les CV, échanges et décisions sont souvent stockés hors d'un référentiel unique.
Il devient difficile de retrouver l'historique d'un candidat, de comparer les
candidatures sur une même offre, ou d'assurer une continuité lorsque plusieurs RH
interviennent.

\paragraph{Manque de traçabilité et de statut clair.}
Sans workflow formalisé (en attente, retenu, non retenu, désisté, admis), le
suivi du pipeline de sélection manque de fiabilité. Les indicateurs de performance
(taux de sélection, délais, volumes) sont alors approximatifs ou absents.

\paragraph{Faible coordination multi-sites.}
Dans un réseau multi-zones, l'absence de périmètre applicatif (droits RH par zone,
supervision admin) peut entraîner des doublons, des incohérences ou une perte de
visibilité managériale.

\paragraph{Expérience candidat dégradée.}
Des délais de réponse longs, l'absence de confirmation d'entretien, ou un parcours
de candidature peu clair nuisent à l'image employeur et augmentent le risque de
désistement.

\paragraph{Peu d'aide à la décision.}
Le tri manuel des CV et l'absence d'outils d'évaluation (QCM, analyse assistée,
synthèses) allongent le temps de présélection et augmentent la subjectivité.

\paragraph{Difficultés de reporting.}
Les exports, tableaux de bord et bilans mensuels restent coûteux à produire lorsque
les données ne sont pas structurées dans un système unique.

Ces limites justifient le recours à une plateforme dédiée, capable d'unifier le
parcours de recrutement tout en respectant l'organisation territoriale de l'entreprise.


\section{Les tendances actuelles des systèmes de recrutement (ATS)}

L'évolution technologique pousse les éditeurs d'ATS et de plateformes RH à innover
afin de répondre aux besoins croissants de performance, de conformité et d'expérience
utilisateur. Les principales tendances observées sont les suivantes~:

\begin{itemize}
  \item \textbf{Centralisation du pipeline candidat}~: une offre unique de vérité pour
        les candidatures, leurs statuts et leur historique.
  \item \textbf{Portails multi-rôles}~: espaces distincts pour l'administration, les RH
        et les candidats, avec gestion fine des habilitations.
  \item \textbf{Automatisation des communications}~: e-mails de convocation, rappels,
        confirmations d'embauche, pièces jointes documentaires.
  \item \textbf{Évaluation assistée}~: QCM de présélection, scoring, analyse de CV
        (extraction de compétences, adéquation au poste), évaluations post-embauche.
  \item \textbf{Pilotage par indicateurs}~: tableaux de bord (volumes, retenus,
        non retenus, désistés, taux de sélection), analyses par agence / poste /
        responsable.
  \item \textbf{Intégration calendrier / visioconférence}~: planification d'entretiens
        en ligne ou physiques, avec restitution dans le dossier candidat.
  \item \textbf{Architecture web moderne}~: applications SPA, API REST, authentification
        sécurisée (souvent token JWT), éventuelle modularité microservices.
  \item \textbf{Conformité et confidentialité}~: protection des données personnelles,
        contrôle d'accès, séparation des responsabilités.
\end{itemize}

Dans ce paysage, une solution \og{}métier\fg{} adaptée à un réseau d'agences ne se
contente pas d'un simple formulaire de candidature~: elle doit refléter la réalité
organisationnelle (zones, agences, RH responsables), industrialiser le suivi et
fournir une base fiable pour le reporting et l'amélioration continue.


\section{Conclusion}

Le recrutement représente aujourd'hui un levier stratégique essentiel dans un
environnement concurrentiel, digitalisé et exigeant en matière d'expérience candidat.
Les approches traditionnelles (outils dispersés, suivi manuel) montrent clairement
leurs limites dès lors que le volume, la multi-sitesité et le besoin de traçabilité
augmentent.

Les systèmes de type ATS s'imposent comme des réponses structurantes aux défis
actuels~: centralisation, automatisation, évaluation, pilotage et gouvernance des
accès. C'est dans cette dynamique que s'inscrit le projet de plateforme intelligente
de gestion du recrutement réalisé au sein de \textbf{DAAM}.

Le chapitre suivant approfondira le \textbf{contexte concret du projet}~: présentation
de l'organisme d'accueil, problématique rencontrée, solution proposée et méthode
de travail retenue (SCRUM).

\clearpage
